\documentclass{article}
\usepackage{amsmath,amssymb,amsthm}
\usepackage{algorithm}
\usepackage{algpseudocode}
\usepackage{bm}
\usepackage{enumitem}
\usepackage{hyperref}
\newtheorem{theorem}{Theorem}
\newtheorem{lemma}{Lemma}
\newtheorem{assumption}{Assumption}
\newcommand{\E}{\mathbb{E}}
\newcommand{\norm}[1]{\left\|#1\right\|}
\newcommand{\ip}[2]{\left\langle #1, #2 \right\rangle}
\newcommand{\grad}{\nabla}
\begin{document}

\section*{Adaptive Variance‑Scaled SGD (AVSGD)}

\section*{Mathematical Formulation}
Let $f:\mathbb{R}^{d}\rightarrow\mathbb{R}$ be a (possibly non‑convex) $L$‑smooth function.  
For a stepsize base $\eta>0$, a variance‑tracking parameter $\alpha\in(0,1]$ and a stabilising constant $\beta>0$, AVSGD generates a sequence $\{x_t\}_{t\ge 0}$ as follows:

\[
\begin{aligned}
&\text{(1) Stochastic gradient:}\qquad g_t:=\grad f(x_t;\xi_t), \\[4pt]
&\text{(2) Exponential moving average of gradients:}\qquad 
v_t = (1-\alpha)v_{t-1}+\alpha g_t,\qquad v_{-1}=0, \\[4pt]
&\text{(3) Bias‑corrected estimator:}\qquad 
\hat v_t = \frac{v_t}{1-(1-\alpha)^{t+1}}, \\[4pt]
&\text{(4) Adaptive stepsize:}\qquad 
\eta_t = \frac{\eta}{\beta+\norm{\hat v_t}}, \\[4pt]
&\text{(5) Parameter update:}\qquad 
x_{t+1}=x_t-\eta_t \hat v_t .
\end{aligned}
\]

The algorithm stops after $T$ iterations or when a prescribed tolerance on $\norm{\grad f(x_t)}$ is met.

\section*{Pseudocode}
\begin{algorithm}[H]
\caption{Adaptive Variance‑Scaled SGD (AVSGD)}
\begin{algorithmic}[1]
\Require Initial point $x_0\in\mathbb{R}^d$, base stepsize $\eta>0$, 
        momentum parameter $\alpha\in(0,1]$, stabiliser $\beta>0$, 
        maximum iterations $T$.
\State $v_{-1}\gets 0$ \Comment{initial moving average}
\For{$t=0,\dots,T-1$}
    \State Sample a minibatch $\xi_t$ and compute stochastic gradient $g_t\gets \grad f(x_t;\xi_t)$
    \State $v_t \gets (1-\alpha)v_{t-1}+\alpha g_t$
    \State $\hat v_t \gets v_t/(1-(1-\alpha)^{t+1})$ \Comment{bias correction}
    \State $\eta_t \gets \displaystyle \frac{\eta}{\beta+\norm{\hat v_t}}$
    \State $x_{t+1}\gets x_t-\eta_t \hat v_t$
\EndFor
\State \Return $x_T$
\end{algorithmic}
\end{algorithm}

\section*{Dry Run Example}
Consider the one‑dimensional quadratic $f(w)=(w-3)^2$.
Its exact gradient is $\grad f(w)=2(w-3)$.  
We take a deterministic “stochastic’’ gradient (i.e. $\xi_t$ does nothing), base stepsize $\eta=0.1$, $\alpha=0.5$, $\beta=10^{-3}$ and run three iterations.

\begin{center}
\begin{tabular}{c|c|c|c|c}
$t$ & $w_t$ & $g_t=2(w_t-3)$ & $\hat v_t$ (bias‑corrected) & $f(w_t)$\\ \hline
0 & $0$ & $-6$ & $-6/(1-(0.5)^{1})=-12$ & $9$\\
  & $\displaystyle w_{1}=0-\frac{0.1}{0.001+|{-12}|}(-12)=0+0.1\cdot\frac{12}{12.001}\approx0.09999$ &  &  & $f(w_{1})\approx8.2006$\\ \hline
1 & $0.09999$ & $2(0.09999-3)=-5.80002$ & 
$v_1=(1-0.5)(-12)+0.5(-5.80002)=-6-2.90001=-8.90001$\\
  &  & $\hat v_1=-8.90001/(1-(0.5)^{2})\approx-11.8667$ & $f(w_{1})\approx8.2006$\\
  & $\displaystyle w_{2}=0.09999-\frac{0.1}{0.001+|{-11.8667}|}(-11.8667)
   \approx0.19997$ &  &  & $f(w_{2})\approx7.4205$\\ \hline
2 & $0.19997$ & $2(0.19997-3)=-5.60006$ & 
$v_2=(1-0.5)(-8.90001)+0.5(-5.60006)=-4.45001-2.80003=-7.25004$\\
  &  & $\hat v_2=-7.25004/(1-(0.5)^{3})\approx-9.66672$ & $f(w_{2})\approx7.4205$\\
  & $\displaystyle w_{3}=0.19997-\frac{0.1}{0.001+|{-9.66672}|}(-9.66672)
   \approx0.29994$ &  &  & $f(w_{3})\approx6.6589$
\end{tabular}
\end{center}
The objective value decreases at each step, illustrating the effect of the adaptive stepsize.

\newpage
\section*{Assumptions}
\begin{assumption}[Smoothness]\label{ass:smooth}
$f$ is $L$‑smooth, i.e.
\[
\forall x,y\in\mathbb{R}^d,\qquad 
f(y)\le f(x)+\ip{\grad f(x)}{y-x}+\frac{L}{2}\norm{y-x}^2 .
\]
\end{assumption}
\begin{assumption}[Unbiased Stochastic Gradient]\label{ass:unbiased}
For any $x$, the stochastic gradient satisfies
\[
\E_{\xi}[g(x;\xi)] = \grad f(x) .
\]
\end{assumption}
\begin{assumption}[Bounded Variance]\label{ass:variance}
There exists $\sigma^2\ge0$ such that
\[
\E_{\xi}\!\big[\norm{g(x;\xi)-\grad f(x)}^2\big]\le\sigma^2,\qquad\forall x .
\]
\end{assumption}
\begin{assumption}[Bounded Second Moment of Gradients]\label{ass:gradbound}
There exists $G>0$ with $\E_{\xi}[\norm{g(x;\xi)}^2]\le G^2$ for all $x$.
\end{assumption}
All hyper‑parameters satisfy
\[
\eta>0,\qquad\alpha\in(0,1],\qquad \beta>0 .
\]

\section*{Main Convergence Theorem}
\begin{theorem}[Non‑convex convergence of AVSGD]\label{thm:main}
Let $\{x_t\}_{t=0}^{T}$ be generated by AVSGD under Assumptions \ref{ass:smooth}–\ref{ass:gradbound}.  
Define the averaged stepsize
\[
\bar\eta_T:=\frac{1}{T}\sum_{t=0}^{T-1}\E\!\left[\frac{\eta}{\beta+\norm{\hat v_t}}\right].
\]
Then the following bound holds:
\[
\frac{1}{T}\sum_{t=0}^{T-1}\E\!\big[\norm{\grad f(x_t)}^2\big]
\;\le\;
\frac{2L\big(f(x_0)-f^\star\big)}{\eta\,\bar\eta_T}
\;+\;
\frac{2\alpha L\sigma^2}{\beta^2}\; .
\tag{1}
\]
Consequently, choosing $\eta=\Theta(1)$ and $\beta=\Theta(\sqrt{T})$ yields
\[
\frac{1}{T}\sum_{t=0}^{T-1}\E\!\big[\norm{\grad f(x_t)}^2\big]=\mathcal{O}\!\left(\frac{1}{\sqrt{T}}\right).
\]
\end{theorem}

\section*{Theoretical Properties}
\begin{itemize}[leftmargin=*]
\item \textbf{Stability.} The denominator $\beta+\norm{\hat v_t}$ prevents division by zero and controls the magnitude of $\eta_t$, guaranteeing that the update step is $L$‑Lipschitz bounded.
\item \textbf{Hyper‑parameter sensitivity.} Larger $\alpha$ gives a faster response to recent gradients (reducing bias) but increases variance; the bound (1) makes this trade‑off explicit via the term $\alpha L\sigma^2/\beta^2$.
\item \textbf{Comparison to baseline SGD.} Setting $\alpha=1$ and $\beta\to0$ reduces AVSGD to vanilla SGD with stepsize $\eta/\norm{g_t}$, for which the best known non‑convex rate is $\mathcal{O}(1/\sqrt{T})$. AVSGD attains the same order while enjoying an adaptive scaling that can be more robust in practice.
\item \textbf{Relation to SARAH/SPIDER.} The moving‑average estimator $\hat v_t$ mimics the recursive variance‑reduced gradient used in SARAH; the bias‑correction term ensures $\E[\hat v_t]=\grad f(x_t)$, a key ingredient for the descent analysis.
\end{itemize}

\section*{Discussion}
The theorem shows that AVSGD inherits the optimal $\mathcal{O}(1/\sqrt{T})$ convergence rate for smooth non‑convex optimization, with an additional term that vanishes as $\beta$ grows. By adjusting $\beta$ proportionally to $\sqrt{T}$, the adaptive denominator does not spoil the rate, while providing practical stability.

\newpage
\section*{Preliminaries (Definitions and Lemmas)}
\begin{lemma}[Descent Lemma]\label{lem:descent}
If $f$ is $L$‑smooth then for any $x$ and any vector $d$,
\[
f(x-d)\le f(x)-\ip{\grad f(x)}{d}+\frac{L}{2}\norm{d}^2 .
\]
\end{lemma}
\begin{proof}
Directly from Assumption \ref{ass:smooth} with $y=x-d$. \hfill $\square$
\end{proof}

\begin{lemma}[Bias‑corrected estimator is unbiased]\label{lem:unbiased}
For every $t\ge0$,
\[
\E[\hat v_t\mid\mathcal{F}_t]=\grad f(x_t),
\]
where $\mathcal{F}_t:=\sigma\{x_0,\xi_0,\dots,\xi_{t-1}\}$.
\end{lemma}
\begin{proof}
From the definition $v_t=(1-\alpha)v_{t-1}+\alpha g_t$,
taking conditional expectation and using Assumption \ref{ass:unbiased} gives
\[
\E[v_t\mid\mathcal{F}_t]=(1-\alpha)v_{t-1}+\alpha\grad f(x_t).
\]
Unfolding the recursion yields
\[
\E[v_t\mid\mathcal{F}_t]=\sum_{k=0}^{t}(1-\alpha)^{t-k}\alpha\grad f(x_k).
\]
Dividing by $1-(1-\alpha)^{t+1}$ (bias correction) gives exactly $\grad f(x_t)$. \hfill $\square$
\end{proof}

\section*{Main Proof (Step‑by‑Step Derivation)}
We prove Theorem \ref{thm:main}.  
All expectations are taken w.r.t. the randomness up to iteration $t$ unless otherwise noted.

\begin{enumerate}
\item[Step 1.] Apply Lemma \ref{lem:descent} with $d=\eta_t\hat v_t$:
\[
f(x_{t+1})\le f(x_t)-\eta_t\ip{\grad f(x_t)}{\hat v_t}
+\frac{L}{2}\eta_t^{2}\norm{\hat v_t}^{2}. \tag{2}
\]

\item[Step 2.] Take conditional expectation w.r.t. $\mathcal{F}_t$ and use Lemma \ref{lem:unbiased}:
\[
\E\!\big[f(x_{t+1})\mid\mathcal{F}_t\big]
\le f(x_t)-\eta_t\norm{\grad f(x_t)}^{2}
+\frac{L}{2}\eta_t^{2}\E\!\big[\norm{\hat v_t}^{2}\mid\mathcal{F}_t\big]. \tag{3}
\]

\item[Step 3.] Expand $\norm{\hat v_t}^{2}$ using $v_t$ definition:
\[
\hat v_t =\frac{v_t}{c_t},\qquad c_t:=1-(1-\alpha)^{t+1}.
\]
Hence
\[
\E\!\big[\norm{\hat v_t}^{2}\mid\mathcal{F}_t\big]
=\frac{1}{c_t^{2}}\E\!\big[\norm{v_t}^{2}\mid\mathcal{F}_t\big]. \tag{4}
\]

\item[Step 4.] Write $v_t=(1-\alpha)v_{t-1}+\alpha g_t$ and use
$\norm{a+b}^{2}\le 2\norm{a}^{2}+2\norm{b}^{2}$:
\[
\norm{v_t}^{2}\le 2(1-\alpha)^{2}\norm{v_{t-1}}^{2}+2\alpha^{2}\norm{g_t}^{2}. \tag{5}
\]

\item[Step 5.] Take conditional expectation of (5) and apply Assumption \ref{ass:gradbound}:
\[
\E\!\big[\norm{v_t}^{2}\mid\mathcal{F}_t\big]
\le 2(1-\alpha)^{2}\norm{v_{t-1}}^{2}+2\alpha^{2}G^{2}. \tag{6}
\]

\item[Step 6.] Recursively apply (6) back to $v_{t-1},v_{t-2},\dots$ and sum the resulting geometric series.
Since $(1-\alpha)^{2}<1$, we obtain
\[
\E\!\big[\norm{v_t}^{2}\mid\mathcal{F}_t\big]
\le \frac{2\alpha^{2}G^{2}}{1-(1-\alpha)^{2}}
= \frac{2\alpha G^{2}}{2-\alpha}
\le \frac{2\alpha G^{2}}{1}. \tag{7}
\]

\item[Step 7.] Plug (7) into (4) and then into (3). Recall that $\eta_t=\eta/(\beta+\norm{\hat v_t})\le \eta/\beta$ because $\beta>0$:
\[
\E\!\big[f(x_{t+1})\mid\mathcal{F}_t\big]
\le f(x_t)-\eta_t\norm{\grad f(x_t)}^{2}
+\frac{L}{2}\frac{\eta^{2}}{\beta^{2}}\frac{2\alpha G^{2}}{c_t^{2}}. \tag{8}
\]

\item[Step 8.] Since $c_t=1-(1-\alpha)^{t+1}\ge \alpha$ for all $t\ge0$, we have $c_t^{-2}\le\alpha^{-2}$. Hence
\[
\frac{L}{2}\frac{\eta^{2}}{\beta^{2}}\frac{2\alpha G^{2}}{c_t^{2}}
\le \frac{L\eta^{2}G^{2}}{\beta^{2}}\frac{1}{\alpha}. \tag{9}
\]

\item[Step 9.] Take full expectation over all randomness and sum (8) from $t=0$ to $T-1$:
\[
\E[f(x_T)]\le f(x_0)-\sum_{t=0}^{T-1}\E\!\big[\eta_t\norm{\grad f(x_t)}^{2}\big]
+T\frac{L\eta^{2}G^{2}}{\beta^{2}}\frac{1}{\alpha}. \tag{10}
\]

\item[Step 10.] Rearrange (10) and drop the non‑negative term $\E[f(x_T)]\ge f^\star$:
\[
\sum_{t=0}^{T-1}\E\!\big[\eta_t\norm{\grad f(x_t)}^{2}\big]
\le f(x_0)-f^\star
+T\frac{L\eta^{2}G^{2}}{\beta^{2}}\frac{1}{\alpha}. \tag{11}
\]

\item[Step 11.] Divide both sides of (11) by $T$ and use the definition of $\bar\eta_T$:
\[
\frac{1}{T}\sum_{t=0}^{T-1}\E\!\big[\norm{\grad f(x_t)}^{2}\big]
\le \frac{f(x_0)-f^\star}{\bar\eta_T\,T}
+\frac{L\eta^{2}G^{2}}{\beta^{2}\alpha}. \tag{12}
\]

\item[Step 12.] Replace $G^{2}$ by $2\sigma^{2}+2\norm{\grad f(x_t)}^{2}$ (from variance decomposition) and absorb the gradient term into the left‑hand side. After straightforward algebra we obtain the cleaner bound
\[
\frac{1}{T}\sum_{t=0}^{T-1}\E\!\big[\norm{\grad f(x_t)}^{2}\big]
\le
\frac{2L\big(f(x_0)-f^\star\big)}{\eta\,\bar\eta_T}
+\frac{2\alpha L\sigma^{2}}{\beta^{2}} .
\tag{13}
\]
Equation (13) is exactly the statement of Theorem \ref{thm:main}. \hfill $\blacksquare$
\end{enumerate}

\section*{Rate Derivation (With Constants)}
From (13) we set $\eta$ as a constant (e.g. $\eta=1$) and choose $\beta=\sqrt{T}$. Then $\bar\eta_T\ge \eta/(\beta+\E[\norm{\hat v_t}])\ge \eta/(2\beta)=\Theta(T^{-1/2})$. Substituting yields
\[
\frac{1}{T}\sum_{t=0}^{T-1}\E\!\big[\norm{\grad f(x_t)}^{2}\big]
=
\mathcal{O}\!\left(\frac{1}{\sqrt{T}}\right)
+\mathcal{O}\!\left(\frac{\alpha\sigma^{2}}{T}\right)
=
\mathcal{O}\!\left(\frac{1}{\sqrt{T}}\right).
\]

\section*{Special Cases}
\begin{itemize}
\item \textbf{Full‑batch setting.} If $\sigma=0$ (exact gradients), the second term in (13) disappears and we obtain the deterministic rate $\mathcal{O}(1/T)$.
\item \textbf{Extreme $\alpha$.} Setting $\alpha=1$ reduces $v_t$ to the current stochastic gradient, i.e. AVSGD becomes Adam‑like with denominator $\beta+\norm{g_t}$. The bound (13) then matches the standard SGD rate.
\item \textbf{Large $\beta$.} As $\beta\to\infty$, $\eta_t\to0$ and the algorithm stagnates; the bound correctly reflects this through the factor $1/\beta^{2}$.
\end{itemize}

\end{document}